% defer/rcurelated.tex
% mainfile: ../perfbook.tex
% SPDX-License-Identifier: CC-BY-SA-3.0

\subsection{RCU 相关工作}
\label{sec:defer:RCU Related Work}
\OriginallyPublished{sec:defer:RCU Related Work}{RCU Related Work}{Linux Weekly News}{PaulEMcKenney2014ReadMostly}
\OriginallyPublished{sec:defer:RCU Related Work}{RCU Related Work}{Linux Weekly News}{PaulEMcKenney2015ReadMostly}

已知最早提及类似 RCU 机制的是
\ppl{Donald}{Knuth}~\cite[page 413 of Fundamental Algorithms]{Knuth73}
针对\ppl{Joseph}{Weizenbaum}的 FORTRAN SLIP 列表处理
工具~\cite{Weizenbaum:1963:SLP:367593.367617}的一个 bug 报告。
Knuth 报告这个 bug 是正当的，因为 SLIP 没有任何
grace-period 保证的概念。

已知最早的非 bug 报告形式的类似 RCU 机制的提及出现在
\pplsur{H. T.}{Kung}和\pplsur{Philip L.}{Lehman}的
里程碑式论文~\cite{Kung80}中。
学术界有一些关于此技术的额外使用~\cite{Manber82,Manber84,BarbaraLiskov1988ArgusCACM,Pugh90,Andrews91textbook,Pu95a,Cowan96a,Rastogi:1997:LPV:645923.671017,Gamsa99}，
但该领域的大量工作反而由
从业者完成~\cite{RichardRashid87a,Hennessy89,Jacobson93,AjuJohn95,Slingwine95,Slingwine97,Slingwine98,McKenney98}。

\QuickQuiz{
	垃圾回收器？
	被动序列化？
	系统参考点？
	Quiescent states？
	老化？
	代？
	为什么这些早期论文的研究人员就不能统一一下术语呢？？？
}\QuickQuizAnswer{
	有多个独立发明的机制大致类似于 RCU\@。
	每组发明者都不知道其他人的存在，
	因此每组理所当然地创造了自己的术语。
	而不同的术语使得任何一组都很难找到其他人。

	抱歉，但生活有时候就是这样！
}\QuickQuizEnd

到 2000 年，主导权已经转移到开源项目，
最显著的是 Linux 内核
社区~\cite{RustyRussell2000a,RustyRussell2000b,McKenney01b,McKenney01a,McKenney02a,Arcangeli03}。\footnote{
	在本书的 {\LaTeX} 源码中的 \co{bib/RCU.bib}
	可以找到一个包含超过 200 个条目的引用列表。}
RCU 于 2002 年末被 Linux 内核接受，
随后在可扩展性、健壮性、实时响应、
能源效率和专用使用场景方面进行了许多改进。
截至 2023 年，Linux 内核 RCU 仍在积极开发中。

\QuickQuiz{
	为什么 Kung 和 Lehman 的论文没有导致 RCU 被立即使用？
}\QuickQuizAnswer{
	一个原因是 Kung 和 Lehman 确实走在了时代前面。
	另一个原因是他们的方法虽然具有开创性，
	但没有考虑到许多软件工程和性能问题。

	要看出他们走在时代前面，考虑一下在他们的论文发表三年后，
	Paul 正在一个运行 BSD 2.8 的 PDP-11 系统上工作。
	这个系统缺乏任何自动配置，
	这意味着任何硬件修改，包括添加新磁盘驱动器，
	都需要手动编辑和重建内核。
	此外，这是一个单 CPU 系统，
	这意味着全系统同步只需简单地禁用中断。

	快进若干年，允许运行时更改硬件配置的多核系统
	变得司空见惯。
	这意味着 1980 年代内核源代码中隐式表示的
	硬件配置数据现在是一个可变的数据结构，
	每次 I/O 都会被访问。
	这样的数据结构很少更改，但可以随时更改。
	而这种读多属性适用于许多其他新时代的数据结构，
	包括关于网络（1980 年代很少见）、
	安全策略（1980 年代用物理锁）、
	软件配置（1980 年代运行时不可变）等等的数据结构。
	因此，1990 年代和 2000 年代比 1980 年代
	有更多机会让 RCU 展示其优势。

	Kung 和 Lehman 的软件工程方面的不足包括
	未标记读者（因此带来调试困难）、
	未提供干净的 RCU API（因此将其机制绑定到特定数据结构），
	以及不允许 grace period 后操作除释放内存外的任何其他操作
	（因此排除了许多 RCU 使用场景）。

	Kung 和 Lehman 提出了两种垃圾回收策略。
	第一种等待在给定时间运行的所有进程终止，
	这代表了另一个软件工程不足，
	排除了其机制在无限期运行的软件中的使用。
	第二种使用每对象引用计数，
	这大大复杂化了其读侧代码（因此代表又一个软件工程不足），
	并且在现代硬件上导致严重的缓存未命中开销
	（因此代表性能不足，参见例如
	\cref{fig:defer:Performance of RCU vs. Reference Counting,fig:defer:Performance of Preemptible RCU vs. Reference Counting}）。

	尽管软件工程和性能方面有这么长的不足清单，
	Kung 和 Lehman 的论文仍然是一项真正令人印象深刻的工作，
	特别是考虑到后来的许多工作
	（无论是独立的还是非独立的）
	都犯了同样的不足，甚至还有其他不足。
}\QuickQuizEnd

然而，在 2010 年代中期，RCU 的研究和开发在许多
社区和机构中出现了可喜的
回升~\cite{FransKaashoek2015ParallelOSHistory}。
\Cref{sec:defer:RCU Uses}描述 RCU 的使用，
\cref{sec:defer:RCU Implementations}描述 RCU 实现
（以及既创建又使用实现的工作），
最后，
\cref{sec:defer:RCU Validation}描述 RCU 及其使用的
验证和确认。

\subsubsection{RCU 的使用}
\label{sec:defer:RCU Uses}

\ppl{Phil}{Howard}和\ppl{Jon}{Walpole}（波特兰州立大学，PSU）
将 RCU 应用于红黑
树~\cite{PhilHowardPhD,PhilHoward2011RCUTMRBTree}，
结合使用软件事务内存同步的更新。
\ppl{Josh}{Triplett}和\ppl{Jon}{Walpole}（同样来自 PSU）
将 RCU 应用于可调整大小的
哈希表~\cite{JoshTriplettPhD,Triplett:2011:RPHash,JonCorbet2014RCUhash1,JonCorbet2014RCUhash2}。
其他 RCU 保护的可调整大小哈希表由
\ppl{Herbert}{Xu}~\cite{HerbertXu2010RCUResizeHash}和
\ppl{Mathieu}{Desnoyers}~\cite{PaulMcKenney2013LWNURCUhash}创建。

\ppl{Austin}{Clements}、\ppl{Frans}{Kaashoek}和\ppl{Nickolai}{Zeldovich}
（MIT）创建了一种 RCU 优化的平衡二叉树
（Bonsai）~\cite{AustinClements2012RCULinux:mmapsem}，
并将此树应用于 Linux 内核的 VM 子系统，
以减少 Linux 内核 \co{mmap_sem} 上的读侧竞争。
这项工作在一个以大量次要缺页错误为特征的微基准测试中
产生了数量级的加速，可扩展到至少 80 个 CPU。
这类似于\ppl{Peter}{Zijlstra}~\cite{PeterZijlstra2014SpeculativePageFault}
较早开发的补丁，两者都受到当时文件系统数据结构
对 RCU 读者不安全这一事实的限制。
\pplsur{Austin}{Clements}等人通过仅优化匿名页面的
缺页路径来避免此限制。
最近，文件系统数据结构已经对 RCU 读者变得
安全~\cite{JonathanCorbet2010dcacheRCU,JonathanCorbet2011dcacheRCUbug}，
因此也许这项工作可以为所有页面类型实现，
而不仅仅是匿名页面——\ppl{Peter}{Zijlstra}实际上最近已经
做了原型，\ppl{Laurent}{Dufour}和\ppl{Michel}{Lespinasse}
继续沿这些方向工作。
\ppl{Matthew}{Wilcox}、\ppl{Liam}{Howlett}和
\ppl{Suren}{Baghdasaryan}正在致力于使用 RCU 来实现对其他
内存管理数据结构的细粒度锁定和无锁访问。

\ppl{Yandong}{Mao}和\ppl{Robert}{Morris}（MIT）以及
\ppl{Eddie}{Kohler}（哈佛大学）创建了另一种名为
Masstree~\cite{Mao:2012:CCF:2168836.2168855}的 RCU 保护树，
结合了 B+ 树和 trie 的思想。
虽然此树比 RCU 保护的哈希表慢约 2.5 倍，
但它支持对键范围的操作，这与哈希表不同。
此外，Masstree 支持具有长共享键前缀的对象的高效存储，
并且还通过日志记录到大容量存储提供持久性。

该论文指出 Masstree 的性能可与 memcached 媲美，
即使 Masstree 持久存储更新而 memcached 没有。
该论文还将 Masstree 的性能与持久数据存储
MongoDB、VoltDB 和 Redis 进行了比较，
报告了 Masstree 的显著性能优势，
在某些情况下超过两个数量级。
另一篇论文~\cite{Tu:2013:STM:2517349.2522713}，
由\ppl{Stephen}{Tu}、\ppl{Wenting}{Zheng}、\ppl{Barbara}{Liskov}
和\ppl{Samuel}{Madden}（MIT）以及\pplsur{Eddie}{Kohler}撰写，
将 Masstree 应用于名为 Silo 的内存数据库，
在一个著名的事务处理基准测试上实现了每秒 700K 事务
（每分钟 42M 事务）。
有趣的是，Silo 在不持有锁的情况下保证线性化，
同时不承担 grace periods 的开销。

\ppl{Maya}{Arbel}和\ppl{Hagit}{Attiya}（以色列理工学院）
采用了更严格的
方法~\cite{MayaArbel2014RCUtree}来实现 RCU 保护的搜索树，
与 Masstree 一样允许并发更新。
该论文包括正确性证明，包括该树上所有操作
都是\IX{linearizable}的证明。
不幸的是，此实现通过在持有锁时承担 grace-period 等待的
全部延迟来实现线性化，这降低了仅更新工作负载的可扩展性。
解决此问题的一种方法是放弃
线性化~\cite{AndreasHaas2012FIFOisnt,PaulEMcKennneyAtomicTreeN4037}，
然而 Arbel 和 Attiya 创建了一种减少低端
grace-period 延迟的 RCU 变体。
当然，没有免费的午餐，这种 RCU 变体似乎在约 32 个 CPU 处
达到可扩展性极限。
虽然放弃线性化以获得性能和可扩展性有很多可取之处，
但看到学者们尝试替代 RCU 实现是非常好的。

\subsubsection{RCU 实现}
\label{sec:defer:RCU Implementations}

\ppl{Timothy}{Harris}创建了一种基于时间的用户空间
RCU~\cite{TimothyLHarris2001}，改进了之前由
Jacobson~\cite{Jacobson93}和 John~\cite{AjuJohn95}创建的版本。
之前这两种基于时间的方法各自假设读者持续时间有严格上限，
这在硬实时系统中可以正确工作。
在非实时系统中，当读者被中断、抢占或以其他方式延迟时，
这种方法容易失败。
然而，这种易失败的实现被独立发明两次的事实
表明了对类似 RCU 机制的深层需求。
\ppl{Timothy}{Harris}通过要求每个读者在开始其
读侧遍历之前获取全局时间基准的快照来改进这两项早期工作。
延迟释放读者可见对象，直到所有进程的读者快照
指示该对象移除之后的时间。
然而，全局时间基准在某些系统上可能既昂贵又不准确。

\ppl{Keir}{Fraser}创建了一种名为\IXacr{ebr}的用户空间 RCU，
用于非阻塞同步和软件事务
内存~\cite{KeirAnthonyFraserPhD,UCAM-CL-TR-579,KeirFraser2007withoutLocks}。
这项工作通过用软件计数器替换全局时钟来改进\ppl{Timothy}{Harris}的工作，
从而消除了当时商品系统全局时钟的大部分开销和所有不准确性。
有趣的是，这项工作一方面引用了 Linux 内核 RCU，
但也为原始的不可抢占 Linux 内核 RCU 实现
启发了\acr{qsbr}的名称。

\ppl{Mathieu}{Desnoyers}创建了用于跟踪的用户空间
RCU~\cite{MathieuDesnoyers2009URCU,MathieuDesnoyersPhD,MathieuDesnoyers2012URCU,PaulMcKenney2013LWNURCU,PaulMcKenney2013LWNURCUhash,PaulMcKenney2013LWNURCUhashAPI,PaulMcKenney2013LWNURCUqueuestack,PaulMcKenney2013LWNURCUqueuestackAPI,PaulMcKenney2013LWNURCUatomicop,PaulMcKenney2013LWNURCUmenagerie,PaulMcKenney2013LWNURCUAPI,PaulMcKenney2013LWNURCUlist,PaulMcKenney2013LWNURCUmenagerieRCU}，
已在许多项目中使用~\cite{MikeDay2013RCUqemu}。

布拉格查理大学的研究人员也在研究 RCU 实现，
包括\ppl{Andrej}{Podzimek}~\cite{AndrejPodzimek2010masters}和
\ppl{Adam}{Hraska}~\cite{AdamHraska2013RCUHelenOS}的论文。

\ppl{Yujie}{Liu}（利哈伊大学）、\ppl{Victor}{Luchangco}（Oracle 实验室）和
\ppl{Michael}{Spear}（同样利哈伊大学）~\cite{Liu:2013:MSA:2549695.2549732}
将可扩展非零指示器
（SNZI）~\cite{FaithEllen:2007:SNZI}用作 grace-period 机制。
其预期用途是实现软件事务内存
（见\cref{sec:future:Transactional Memory}），
这施加了线性化要求，反过来似乎限制了可扩展性。

类似 RCU 的机制也在进入 Java。
\pplsur{KC}{Sivaramakrishnan}等人~\cite{Sivaramakrishnan:2012:ERB:2258996.2259005}
使用类似 RCU 的机制消除了与 Java 垃圾回收器
交互时所需的读屏障，从而显著提高了性能。

\ppl{Ran}{Liu}、\ppl{Heng}{Zhang}和\ppl{Haibo}{Chen}（上海交通大学）
创建了 RCU 的专用变体，用于优化的
``被动读写锁''~\cite{RanLiu2014PassiveRWLock}，
类似于\ppl{Gautham}{Shenoy}~\cite{GauthamShenoy2006RCUrwlock}和
\ppl{Srivatsa}{Bhat}~\cite{SrivatsaSBhat2014RCUrwlock}创建的版本。
Liu 等人的论文从多个角度来看都很有趣~\cite{PaulEMcKenney2014ReadMostly}。

\ppl{Mike}{Ash}发布了~\cite{MikeAsh2015Apple}Apple Objective-C
运行时中类似 RCU 原语的描述。
此方法通过指定的代码范围标识读侧临界区，
因此有资格作为另一种实现零读侧开销的方法，
虽然对跨多个函数的大型读侧临界区提出了一些有趣的实际挑战。

\ppl{Pedro}{Ramalhete}和\ppl{Andreia}{Correia}~\cite{PedroRmalhete2015PoorMansRCU}
产生了``穷人的 RCU''，尽管使用了一对
\IXhpl{reader-writer}{lock}，仍然为
读者提供了\IXalt{lock-free}{lock free}的\IXpl{forward-progress guarantee}~\cite{PaulEMcKenney2015ReadMostly}。

\ppl{Maya}{Arbel}和\ppl{Adam}{Morrison}~\cite{Arbel:2015:PRR:2858788.2688518}
产生了``谓词 RCU''，它努力缩短 grace-period 持续时间，
以高效支持在 grace periods 期间持有更新侧锁的算法。
这导致减少了更新批处理到 grace periods 中的数量和可扩展性，
但确实成功提供了短 grace periods。

\QuickQuiz{
	为什么不在等待 grace period 之前释放锁，
	或者使用像 \co{call_rcu()} 这样的东西代替等待 grace period？
}\QuickQuizAnswer{
	作者希望支持\IX{linearizable}的树操作，
	使得对树的并发添加、删除和搜索看起来
	以某种全局一致的顺序执行。
	在他们的搜索树中，这需要在 grace periods 期间持有锁。
	（在大多数情况下最好放弃线性化作为要求，
	但线性化是一个出人意料地流行的（且昂贵的！\@）要求。）
}\QuickQuizEnd

\ppl{Alexander}{Matveev}（MIT）、\ppl{Nir}{Shavit}（MIT 和特拉维夫大学）、
\ppl{Pascal}{Felber}（纳沙泰尔大学）和\ppl{Patrick}{Marlier}（同样纳沙泰尔大学）~\cite{Matveev:2015:RLS:2815400.2815406}
产生了一种类似 RCU 的机制，可以被认为是
显式标记只读事务的软件事务内存。
他们的使用场景要求在 grace periods 期间持有锁，
这限制了可扩展性~\cite{PaulEMcKenney2015ReadMostly,PaulEMcKenney2015ReadMostlySidebar}。
这似乎是第一个很好地利用 \co{rcutorture} 测试套件的
学术 RCU 相关工作，也是第一个向 Linux 内核 RCU
提交性能改进并被接受到 v4.4 的工作。

\ppl{Alexander}{Matveev}的 RLU 之后是
\ppl{Jaeho}{Kim}等人的 MV-RLU~\cite{Kim:2019:MSR:3297858.3304040}。
这项工作通过允许多个并发更新、避免在 grace periods 期间
持有锁以及使用异步 grace periods
（例如 \co{call_rcu()} 代替 \co{synchronize_rcu()}）
来改善 RLU 的可扩展性。
该论文还做了一些有趣的性能评估选择，
在\cpageref{sec:future:Deferred Reclamation}上的
\cref{sec:future:Deferred Reclamation}中有进一步讨论。

\ppl{Adam}{Belay}等人创建了一个 RCU 实现，
保护其 IX 操作系统中 TCP/IP 地址解析协议（ARP）
使用的数据结构~\cite{Belay:2016:IOS:3014162.2997641}。

\ppl{Geoff}{Romer}和\ppl{Andrew}{Hunter}（均在 Google）
为将基于 cell 的 API 用于单例数据结构的 RCU 保护
提出了纳入 C++ 标准的
建议~\cite{GeoffRomer2018C++DeferredReclamationP0561R4}。

\ppl{Dimitrios}{Siakavaras}等人将
HTM 和 RCU 应用于搜索树~\cite{Siakavaras2017CombiningHA,DimitriosSiakavaras2020RCU-HTM-B+Trees}，
\ppl{Christina}{Giannoula}等人使用 HTM 和 RCU 进行
图着色~\cite{ChristinaGiannoula2018HTM-RCU-graphcoloring}，
\ppl{SeongJae}{Park}等人使用 HTM 和 RCU 来优化
\IXacr{numa}系统上的高竞争锁。

\ppl{Alex}{Kogan}等人将 RCU 应用于构建可扩展
地址空间的范围锁~\cite{AlexKogan2020RCUrangelocks}。

2023 年 6 月 17 日，ISO C++ 标准委员会投票将 RCU
纳入 C++26~\cite{PaulEMcKennney2023C++26RCU4}。

RCU 的生产使用列在\cref{sec:defer:Production Uses of RCU}中。

\subsubsection{RCU 验证}
\label{sec:defer:RCU Validation}

在 2017 年初，普遍认识到几乎任何 bug 都是潜在的安全漏洞，
因此验证和确认是一等关注事项。

石溪大学的研究人员产生了一个 RCU 感知的数据竞争
检测器~\cite{AbhinavDuggal2010Masters,JustinSeyster2012PhD,Seyster:2011:RFA:2075416.2075425}。
\ppl{Alexey}{Gotsman}（IMDEA）、\ppl{Noam}{Rinetzky}（特拉维夫大学）
和\ppl{Hongseok}{Yang}（牛津大学）发表了
一篇论文~\cite{AlexeyGotsman2012VerifyGraceExtended}，
用分离逻辑表达 RCU 的形式语义，并继续研究并发的其他方面。

\ppl{Joseph}{Tassarotti}（卡内基梅隆大学）、\ppl{Derek}{Dreyer}（马普
软件系统研究所）和\ppl{Viktor}{Vafeiadis}
（同样 MPI-SWS）~\cite{JosephTassarotti2015RCUproof}
产生了用户空间 RCU 的\IXacrf{qsbr}变体的
手动形式正确性证明~\cite{MathieuDesnoyers2009URCU,MathieuDesnoyers2012URCU}。
\ppl{Lihao}{Liang}（牛津大学）、\pplmdl{Paul E.}{McKenney}（IBM）、
\ppl{Daniel}{Kroening}和\ppl{Tom}{Melham}
（均为牛津大学）~\cite{LihaoLiang2016VerifyTreeRCU}
使用\IXacrf{cbmc}~\cite{EdmundClarke2004CBMC}
产生了 Linux 内核 Tree RCU 重要部分的机械正确性证明。
\ppl{Lance}{Roy}~\cite{LanceRoy2017CBMC-SRCU}使用 CBMC 产生了
Linux 内核 sleepable RCU（SRCU）~\cite{PaulEMcKenney2006c}
重要部分的类似正确性证明。
最后，\ppl{Michalis}{Kokologiannakis}和\ppl{Konstantinos}{Sagonas}
（雅典国家技术大学）~\cite{MichalisKokologiannakis2017NidhuggRCU,MichalisKokologiannakis2019RCUstatelessModelCheck}
使用 Nidhugg 工具~\cite{CarlLeonardsson2014Nidhugg}
产生了 Linux 内核 Tree RCU 更大部分的机械正确性证明。

这些工作都没有发现除了专门注入以测试验证工具的 bug 之外的任何 bug。
相比之下，\ppl{Alex}{Groce}（俄勒冈州立大学）、\ppl{Iftekhar}{Ahmed}、
\ppl{Carlos}{Jensen}（均为 OSU）和\pplmdl{Paul E.}{McKenney}
（IBM）~\cite{Groce:2015:VMC:2916135.2916190}
自动变异 Linux 内核 RCU 的源代码以测试
\co{rcutorture} 测试套件的覆盖率。
该工作发现了此套件覆盖率中的几个漏洞，
其中一个隐藏了 Tiny RCU 中的一个真实 bug（已修复）。

如果幸运的话，所有这些验证工作最终将产生
更多更好的工具来验证并发代码。
